\begin{proofbox}
\textbf{预备: 两个焦半径公式从哪里来?}

设 $c^2=a^2-b^2$, $e=c/a$, $F_1(-c,0),F_2(c,0)$.
由于 $P(x_0,y_0)$ 在椭圆上,
\[
y_0^2=b^2\left(1-\frac{x_0^2}{a^2}\right).
\]
于是
\[
\begin{aligned}
r_1^2&=(x_0+c)^2+y_0^2\\
&=x_0^2+2cx_0+c^2+b^2-\frac{b^2x_0^2}{a^2}\\
&=a^2+2cx_0+\frac{c^2x_0^2}{a^2}
=\left(a+\frac{cx_0}{a}\right)^2.
\end{aligned}
\]
因为 $|x_0|\le a$ 且 $a>c$, 故括号内恒为正,
\[
\boxed{r_1=a+ex_0},\qquad \boxed{r_2=a-ex_0},
\qquad\boxed{r_1r_2=a^2-e^2x_0^2}.
\]

\textbf{方法一: 最短代数推导 (考试首选)}

切线的方程和中心到切线的距离分别是
\[
l:\ \frac{x_0x}{a^2}+\frac{y_0y}{b^2}=1,
\qquad
\frac1{d^2}=\frac{x_0^2}{a^4}+\frac{y_0^2}{b^4}.
\]
\textbf{关键只做一次椭圆代换}: $y_0^2/b^2=1-x_0^2/a^2$, 所以
\[
\begin{aligned}
\frac{a^2b^2}{d^2}
&=\frac{b^2x_0^2}{a^2}+\frac{a^2y_0^2}{b^2}\\
&=\frac{b^2x_0^2}{a^2}+a^2-x_0^2\\
&=a^2-\frac{(a^2-b^2)x_0^2}{a^2}\\
&=a^2-e^2x_0^2=r_1r_2.
\end{aligned}
\]
于是 $d^2r_1r_2=a^2b^2$.
因 $a,b,d,r_1,r_2$ 均为正数, 两边取正平方根:
\[
\boxed{\sqrt{r_1r_2}\,d=ab}.
\]
\textbf{记忆点}: ``写切线 $\to$ 距离取倒数平方 $\to$ 椭圆代换 $\to$ 认出焦半径乘积''.

\medskip
\textbf{方法二: 纯几何 (反射性质 + 焦点三角形余弦定理)}

记 $\varphi=\angle F_1PF_2$, $\theta$ 为 $PF_1$ 与切线所成的锐角.
椭圆切线的反射性质保证 $PF_2$ 与切线所成的锐角也为 $\theta$,
并且 $\theta=90^\circ-\varphi/2$.

为说明这一反射性质为何成立: 设 $F_1'$ 是 $F_1$ 关于切线的对称点.
椭圆切线上的任意点 $Q$ 都满足
$QF_1+QF_2\ge PF_1+PF_2=2a$, 等号仅在 $Q=P$.
又因 $QF_1=QF_1'$, 由三角形两边之和不小于第三边可知,
直线上的最短折线路径 $F_1'QF_2$ 在 $F_1',Q,F_2$ 共线时取得.
故 $F_1',P,F_2$ 共线, 两条焦半径与切线所成的角相等.

设 $h_1,h_2$ 分别是焦点 $F_1,F_2$ 到切线的垂距.
由直角三角形 $h_i=r_i\sin\theta$.
两焦点与中心同在切线内侧, 且 $O$ 为 $F_1F_2$ 中点, 故
\[
\begin{aligned}
d&=\frac{h_1+h_2}{2}
=\frac{r_1+r_2}{2}\sin\theta\\
&=a\cos\frac\varphi2.
\end{aligned}
\]
另一方面, 在 $\triangle F_1PF_2$ 中, 余弦定理给出
\[
(2c)^2=r_1^2+r_2^2-2r_1r_2\cos\varphi.
\]
整理为
\[
\begin{aligned}
4r_1r_2\cos^2\frac\varphi2
&=2r_1r_2(1+\cos\varphi)\\
&=(r_1+r_2)^2-4c^2\\
&=4a^2-4c^2=4b^2.
\end{aligned}
\]
所以 $\sqrt{r_1r_2}\cos(\varphi/2)=b$.
与 $d=a\cos(\varphi/2)$ 联立, 得主公式.
整个证明只用了\textbf{切线的反射性质、中点投影、余弦定理和椭圆定义}.

\medskip
\textbf{方法三: 纯几何的另一种算法 (正交投影 + 勾股定理)}

仍设焦半径与切线的锐角为 $\theta$.
沿切线方向和垂线方向分别作投影.
因为 $PF_1,PF_2$ 的切向投影方向相反,
$F_1F_2$ 的\textbf{切向投影长度}为 $(r_1+r_2)\cos\theta$;
而\textbf{法向投影的长度}为 $|r_1-r_2|\sin\theta$.
根据勾股定理,
\[
\begin{aligned}
(2c)^2
&=(r_1+r_2)^2\cos^2\theta
 +(r_1-r_2)^2\sin^2\theta\\
&=(r_1+r_2)^2
 -\bigl[(r_1+r_2)^2-(r_1-r_2)^2\bigr]\sin^2\theta\\
&=4a^2-4r_1r_2\sin^2\theta.
\end{aligned}
\]
因此
\[
r_1r_2\sin^2\theta=a^2-c^2=b^2.
\]
而由中心是焦点中点,
\[
d=\frac{r_1\sin\theta+r_2\sin\theta}{2}
=a\sin\theta.
\]
故 $\sqrt{r_1r_2}\,d
=\bigl(\sqrt{r_1r_2}\sin\theta\bigr)a=ab$.
这个证明不需要坐标, 连余弦定理也不需要, 只需要\textbf{投影和勾股}.

\medskip
\textbf{方法四: 常规解析法 (切线距离公式逐项计算)}

切线化为一般式:
\[
b^2x_0x+a^2y_0y-a^2b^2=0.
\]
中心到切线的距离为
\[
d^2=\frac{a^4b^4}{b^4x_0^2+a^4y_0^2}.
\]
利用 $y_0^2=b^2(1-x_0^2/a^2)$,
\[
\begin{aligned}
b^4x_0^2+a^4y_0^2
&=b^4x_0^2+a^4b^2-a^2b^2x_0^2\\
&=b^2\left[a^4-(a^2-b^2)x_0^2\right]\\
&=a^2b^2\left(a^2-\frac{c^2x_0^2}{a^2}\right)\\
&=a^2b^2r_1r_2.
\end{aligned}
\]
代回即可得 $d^2=a^2b^2/(r_1r_2)$, 从而证得结论.
此法通用稳妥, 但运算量多于方法一.

\medskip
\textbf{方法五: 参数法 (把椭圆看成圆的不等比拉伸)}

将椭圆上的点表示为 $P(a\cos t,b\sin t)$, 则切线为
\[
\frac{x\cos t}{a}+\frac{y\sin t}{b}=1.
\]
点到直线距离给出
\[
\begin{aligned}
d&=\frac1{\sqrt{\cos^2t/a^2+\sin^2t/b^2}}\\
&=\frac{ab}{\sqrt{b^2\cos^2t+a^2\sin^2t}}.
\end{aligned}
\]
而由焦半径公式
\[
\begin{aligned}
r_1r_2&=(a+c\cos t)(a-c\cos t)\\
&=a^2-c^2\cos^2t\\
&=b^2\cos^2t+a^2\sin^2t.
\end{aligned}
\]
两式中的根号内是同一个量, 故 $d\sqrt{r_1r_2}=ab$.
\textbf{数形联系}: 圆对应 $a=b$ 时, 根号内恒为 $a^2$;
椭圆因不等比拉伸而使距离变化, 但两处变化恰好抵消.

\medskip
\textbf{推论: 两个焦点到切线的垂距之积恒为 $b^2$}

沿用方法二、三的 $h_1,h_2,\theta$, 已知
\[
h_1=r_1\sin\theta,\qquad h_2=r_2\sin\theta,
\qquad r_1r_2\sin^2\theta=b^2.
\]
故
\[
\boxed{h_1h_2=b^2},\qquad
\boxed{d=\frac{h_1+h_2}{2}},\qquad
\boxed{\frac{h_1}{h_2}=\frac{r_1}{r_2}}.
\]
进一步验证
\[
\begin{aligned}
d^2&=\frac{(h_1+h_2)^2}{4}
=\frac{(r_1+r_2)^2\sin^2\theta}{4}\\
&=\frac{a^2h_1h_2}{r_1r_2}
=\frac{a^2b^2}{r_1r_2}.
\end{aligned}
\]
这提供了从\textbf{焦点到切线的垂距乘积}回到 C012 的另一条证明链.
\end{proofbox}
